

\begin{unnumbered}

\section{ВИСНОВКИ}

У ході виконання дипломного проєкту розроблено кросплатформну клієнт-серверну систему агрегації та обробки медико-біологічних даних з використанням хмарних технологій - мобільний застосунок для відстеження розвитку немовляти на платформах iOS та Android.

Проведено аналіз предметної області та огляд трьох популярних аналогів (Sprout Baby Tracker, Huckleberry, Glow Baby). Встановлено, що жоден із них не поєднує в собі повноцінного антропометричного моніторингу відповідно до стандартів \gls{ВООЗ}, синхронізації між пристроями в реальному часі.

Обрано технологічний стек: React Native та Expo для клієнтської частини, Supabase з PostgreSQL - для хмарного бекенду. Спроектовано реляційну схему бази даних із захистом на рівні рядків (\gls{RLS}), що забезпечує ізоляцію даних між користувачами на рівні \gls{СУБД}.

Реалізовано такі модулі системи: журнал активностей немовляти (годування, сон, підгузники), моніторинг антропометричних показників із кривими \gls{ВООЗ}, управління вакцинаціями та досягненнями, а також синхронізація між пристроями в реальному часі.

Проведено тестування на рівнях юніт-тестів та функціонального тестування. Застосунок розгорнуто на хмарній інфраструктурі та готовий до використання.

\end{unnumbered}
